\documentclass[11pt,a4paper]{article}

\usepackage{iftex}
\ifPDFTeX
  \usepackage[T1]{fontenc}
  \usepackage[utf8]{inputenc}
  \usepackage{lmodern}
\else
  \usepackage{fontspec}
\fi

\usepackage[brazilian]{babel}
\usepackage[margin=2.5cm]{geometry}
\usepackage{amsmath,amssymb}
\usepackage{array}
\usepackage{booktabs}
\usepackage{enumitem}
\usepackage{hyperref}

\hypersetup{
  colorlinks=true,
  linkcolor=black,
  urlcolor=blue,
  pdftitle={Algoritmo genético para classificação},
  pdfauthor={Fellipe Almassy}
}

\newcommand{\cod}[1]{\texttt{#1}}
\newcommand{\cromo}{\mathbf{c}}

\title{Algoritmo genético para classificação}
\author{Fellipe Almassy}
\date{Novembro de 2026}

\begin{document}

\maketitle

\begin{abstract}
Este documento explica como um algoritmo genético pode treinar um classificador
linear. Em vez de usar gradiente, o método evolui uma população de soluções
candidatas: as melhores têm mais chance de se reproduzir, seus descendentes sofrem
mutações e substituem as piores soluções. O texto descreve cada etapa do algoritmo,
com a notação matemática e um exemplo numérico, e indica a função correspondente da
biblioteca \cod{algoritmo-genetico}.
\end{abstract}

\section{Objetivo}

Partimos de uma base de dados rotulada, formada por lições (amostras). Cada lição tem
um vetor de \emph{features} e uma classe conhecida. O objetivo do algoritmo genético é
encontrar os coeficientes de uma função linear que, a partir das features, indique a
classe de cada lição. Depois do treinamento, essa função pode classificar dados novos.

\section{Notação}

\begin{center}
\begin{tabular}{cl}
\toprule
Símbolo & Significado \\
\midrule
$N$ & número de lições da base \\
$n$ & número de features \\
$K$ & número de classes, numeradas de $0$ a $K-1$ \\
$\mathbf{x} = (x_1, \dots, x_n)$ & features de uma lição \\
$y$ & classe real de uma lição \\
$m$ & número de cromossomos da população \\
$L = n + 1$ & número de genes de cada cromossomo \\
$h$ & número de filhos gerados por geração \\
$r$ & número de genes alterados em cada filho \\
$s$ & número de cromossomos substituídos por geração \\
\bottomrule
\end{tabular}
\end{center}

\section{Representação da solução}

Cada \textbf{cromossomo} é um vetor de números reais que representa uma solução
candidata. Ele é formado por um \textbf{bias} $b$ e por um \textbf{gene} $g_j$ para
cada feature:
\begin{equation}
  \cromo = (b, g_1, g_2, \dots, g_n).
\end{equation}

Com esses valores, o cromossomo calcula uma pontuação $q$ para qualquer lição:
\begin{equation}
  q(\mathbf{x}) = b + g_1 x_1 + g_2 x_2 + \dots + g_n x_n
                = b + \sum_{j=1}^{n} g_j x_j.
  \label{eq:pontuacao}
\end{equation}

Cada gene indica o peso de uma feature na decisão. O bias desloca a pontuação de todas
as lições pelo mesmo valor. Sem ele, uma lição com todas as features iguais a zero
teria sempre $q = 0$, e a fronteira de decisão seria obrigada a passar pela origem.

Na classificação binária, a pontuação é convertida em classe pelo sinal:
\begin{equation}
  \hat{y} =
  \begin{cases}
    1, & \text{se } q(\mathbf{x}) \geq 0,\\
    0, & \text{se } q(\mathbf{x}) < 0.
  \end{cases}
  \label{eq:binario}
\end{equation}

Na classificação ordinal, as $K$ classes têm uma ordem (por exemplo,
ruim $<$ médio $<$ bom). Definimos $K-1$ pontos de corte fixos
$t_1 < t_2 < \dots < t_{K-1}$, escolhidos pelo usuário e que não evoluem durante o
treino. Com as convenções $t_0 = -\infty$ e $t_K = +\infty$, a classe prevista é
\begin{equation}
  \hat{y} = k \quad \text{tal que} \quad t_k \leq q(\mathbf{x}) < t_{k+1}.
  \label{eq:ordinal}
\end{equation}

\section{O ciclo do algoritmo}

\begin{enumerate}[nosep]
  \item Criar a população inicial.
  \item Calcular o fitness de cada cromossomo.
  \item Selecionar um pai e uma mãe por roleta.
  \item Gerar $h$ filhos por cruzamento.
  \item Aplicar mutação aos filhos.
  \item Substituir os $s$ piores cromossomos pelos $s$ melhores filhos e voltar à etapa~2.
\end{enumerate}

\medskip
O ciclo termina quando o melhor fitness encontrado atinge o valor alvo ou quando acaba
o número máximo de gerações. O resultado é o melhor cromossomo visto em todo o treino,
guardado à parte, e não necessariamente um cromossomo da última população.

\subsection*{Etapa 1: criação da população}

A população inicial é uma matriz $P$ com $m$ linhas e $L$ colunas. Cada linha é um
cromossomo e cada coluna é um gene, sendo a primeira coluna a do bias:
\begin{equation}
  P =
  \begin{pmatrix}
    b_1    & g_{1,1} & g_{1,2} & \cdots & g_{1,n} \\
    b_2    & g_{2,1} & g_{2,2} & \cdots & g_{2,n} \\
    \vdots & \vdots  & \vdots  & \ddots & \vdots  \\
    b_m    & g_{m,1} & g_{m,2} & \cdots & g_{m,n}
  \end{pmatrix}.
\end{equation}

Todos os valores são sorteados de maneira uniforme em um intervalo, por padrão
$[-1, 1]$. Os limites são parâmetros da função e podem ser alterados.

\subsection*{Etapa 2: cálculo do fitness}

O fitness é a nota de um cromossomo: quanto melhor ele classifica as lições da base,
maior a nota. A medida mais direta seria a acurácia, isto é, a fração de acertos. Ela,
porém, é enganosa quando as classes são desbalanceadas. Em uma base com 95 lições da
classe 1 e 5 da classe 0, um cromossomo que responde sempre 1 tem 95\% de acurácia sem
ter aprendido nada.

Por isso, medimos o recall de cada classe, ou seja, a fração de lições dessa classe que
o cromossomo acertou:
\begin{equation}
  R_k = \frac{\text{acertos na classe } k}{\text{total de lições da classe } k}.
\end{equation}

O fitness é o produto dos recalls:
\begin{equation}
  F = \prod_{k=0}^{K-1} R_k.
  \label{eq:fitness}
\end{equation}

Assim, $0 \leq F \leq 1$, e $F = 0$ sempre que o cromossomo erra todas as lições de
alguma classe. No caso acima, o cromossomo que responde sempre 1 tem $R_1 = 1$ e
$R_0 = 0$, logo $F = 0$. O resultado depende apenas da fração de acertos dentro de cada
classe, e não de quantas lições cada uma tem.

\paragraph{Exemplo.} Considere a base da Tabela~\ref{tab:base} e o cromossomo
$\cromo = (2,\, 3,\, -4,\, 1,\, -1)$. Os valores são inteiros e estão fora do intervalo
padrão apenas para facilitar as contas. A pontuação fica
\begin{equation*}
  q(\mathbf{x}) = 2 + 3x_1 - 4x_2 + 1x_3 - 1x_4.
\end{equation*}

\begin{table}[ht]
\centering
\caption{Base de dados do exemplo.}
\label{tab:base}
\begin{tabular}{lccccc}
\toprule
 & $x_1$ & $x_2$ & $x_3$ & $x_4$ & Classe \\
\midrule
Lição 1 & 1 & 0 & 1 & 1 & 1 \\
Lição 2 & 1 & 1 & 0 & 0 & 1 \\
Lição 3 & 1 & 0 & 0 & 1 & 0 \\
Lição 4 & 0 & 1 & 1 & 1 & 0 \\
\bottomrule
\end{tabular}
\end{table}

A Tabela~\ref{tab:avaliacao} aplica a Equação~\eqref{eq:pontuacao} a cada lição e
converte a pontuação em classe pela Equação~\eqref{eq:binario}.

\begin{table}[ht]
\centering
\caption{Avaliação do cromossomo em cada lição.}
\label{tab:avaliacao}
\begin{tabular}{llcccc}
\toprule
 & Cálculo & $q$ & Prevista & Real & Resultado \\
\midrule
Lição 1 & $2 + 3\cdot1 - 4\cdot0 + 1\cdot1 - 1\cdot1$ & $5$  & 1 & 1 & acerto \\
Lição 2 & $2 + 3\cdot1 - 4\cdot1 + 1\cdot0 - 1\cdot0$ & $1$  & 1 & 1 & acerto \\
Lição 3 & $2 + 3\cdot1 - 4\cdot0 + 1\cdot0 - 1\cdot1$ & $4$  & 1 & 0 & erro \\
Lição 4 & $2 + 3\cdot0 - 4\cdot1 + 1\cdot1 - 1\cdot1$ & $-2$ & 0 & 0 & acerto \\
\bottomrule
\end{tabular}
\end{table}

O cromossomo acertou 3 das 4 lições, uma acurácia de 75\%. Por classe, temos
\begin{equation*}
  R_1 = \frac{2}{2} = 1, \qquad R_0 = \frac{1}{2} = 0{,}5, \qquad
  F = R_1 \cdot R_0 = 0{,}5.
\end{equation*}

\subsection*{Etapa 3: seleção por roleta}

Os fitnesses são normalizados em probabilidades que somam 1:
\begin{equation}
  p_i = \frac{F_i}{\sum_{l=1}^{m} F_l}.
\end{equation}
Se todos os fitnesses forem zero, todos os cromossomos recebem a mesma chance,
$p_i = 1/m$.

Para girar a roleta, sorteamos $u$ uniformemente em $[0, 1)$ e escolhemos o primeiro
cromossomo cuja probabilidade acumulada ultrapassa $u$. O sorteio é feito duas vezes,
uma para o pai e outra para a mãe, de modo que os dois podem ser o mesmo cromossomo. A
cada geração é escolhido um único casal.

\paragraph{Exemplo.} Com $F = (0{,}5;\ 0{,}25;\ 0{,}25)$, temos
$p = (0{,}5;\ 0{,}25;\ 0{,}25)$, e as probabilidades acumuladas são
$(0{,}5;\ 0{,}75;\ 1)$. Se $u = 0{,}6$, o sorteado é o segundo cromossomo.

\subsection*{Etapa 4: cruzamento de um ponto}

Para cada um dos $h$ filhos, sorteamos um ponto de corte $c \in \{1, \dots, L-1\}$. O
filho recebe os $c$ primeiros genes do pai e os demais genes da mãe:
\begin{equation}
  \cromo^{\text{filho}} =
  (\cromo^{\text{pai}}_{1}, \dots, \cromo^{\text{pai}}_{c},\;
   \cromo^{\text{mãe}}_{c+1}, \dots, \cromo^{\text{mãe}}_{L}).
\end{equation}
Como o corte nunca fica nas pontas, o filho sempre herda ao menos um gene de cada um.

\paragraph{Exemplo.} Com pai $(2,\, 3,\, -4,\, 1,\, -1)$, mãe $(-1,\, 0{,}5,\, 2,\, -3,\, 1)$
e $c = 2$, o filho é $(2,\, 3,\, 2,\, -3,\, 1)$.

\subsection*{Etapa 5: mutação}

Em cada filho, sorteamos $r$ posições diferentes e trocamos o gene de cada uma por um
novo valor uniforme em $[-1, 1]$. A mutação traz para a população valores que não
existiam nos pais e, por isso, ajuda a manter a diversidade. Na implementação atual,
esse intervalo é fixo e não acompanha os limites escolhidos na etapa~1.

\paragraph{Exemplo.} Se a posição sorteada for a do quarto gene e o novo valor for
$0{,}4$, o filho $(2,\, 3,\, 2,\, -3,\, 1)$ passa a ser $(2,\, 3,\, 2,\, 0{,}4,\, 1)$.

\subsection*{Etapa 6: atualização da população}

Calculamos o fitness dos filhos mutados. Os $s$ cromossomos de menor fitness da
população saem e os $s$ filhos de maior fitness entram em seu lugar, com
$0 \leq s \leq \min(m, h)$. O restante da população passa intacto para a próxima
geração, e os demais filhos são descartados.

A troca não compara filhos e pais: um filho pode entrar mesmo com fitness menor que o
do cromossomo que substituiu. Quando $s < m$, os melhores cromossomos da população são
preservados, já que apenas os piores saem.

\section{Parâmetros}

A Tabela~\ref{tab:parametros} relaciona os símbolos deste texto aos parâmetros da
função \cod{treinar}.

\begin{table}[ht]
\centering
\caption{Parâmetros de \cod{treinar} e seus valores padrão.}
\label{tab:parametros}
\begin{tabular}{clc}
\toprule
Símbolo & Parâmetro & Padrão \\
\midrule
$m$ & \cod{qtd\_cromossomos} & 30 \\
$r$ & \cod{qtd\_genes\_mutados} & 2 \\
$h$ & \cod{qtd\_filhos} & 16 \\
$s$ & \cod{qtd\_substituicoes} & metade de $h$ (8) \\
--- & \cod{geracoes} & 1000 \\
--- & \cod{fitness\_alvo} & 1,0 \\
$t_1, \dots, t_{K-1}$ & \cod{thresholds} & nenhum (binário) \\
--- & \cod{minimo}, \cod{maximo} & $-1$ e $1$ \\
\bottomrule
\end{tabular}
\end{table}

\section{Correspondência com o código}

\begin{table}[ht]
\centering
\caption{Onde cada etapa está implementada, em \cod{src/algoritmo\_genetico/}.}
\begin{tabular}{lll}
\toprule
Etapa & Arquivo & Função \\
\midrule
1. Criação    & \cod{criacao.py}    & \cod{criar\_cromossomos} \\
2. Fitness    & \cod{fitness.py}    & \cod{fitness\_binario}, \cod{fitness\_ordinal} \\
3. Seleção    & \cod{fitness.py}    & \cod{fitness\_percentual} \\
              & \cod{selecao.py}    & \cod{selecionar\_pais\_roleta} \\
4. Cruzamento & \cod{cruzamento.py} & \cod{cruzar\_um\_ponto} \\
5. Mutação    & \cod{mutacao.py}    & \cod{mutar} \\
6. Atualização & \cod{atualizacao.py} & \cod{atualizar\_populacao} \\
Ciclo completo & \cod{nucleo.py}    & \cod{treinar} \\
Classificação & \cod{nucleo.py}     & \cod{prever} \\
\bottomrule
\end{tabular}
\end{table}

\end{document}
